% ===================================================================
%  TAFEL Nr. 12 — Der Bahnhof. --- Die Eisenbahn. --- Die Schiffe.
%  Traduction 2026 d'apres texto/io, controlee sur texto/fr et
%  texto/en. Consigne : voir texto/de/10-tafel-01.tex.
%
%  LE VOYAGE DE CETTE PAGE VA EN ALLEMAGNE, ET LE VOYAGEUR EST
%  FRANCAIS. C'est la seule des seize planches ou la colonne allemande
%  traduise un texte qui parle d'elle, et de l'autre cote de la
%  frontiere. On ne change rien : Guignon fait descendre son
%  narrateur a « Deutsch-Avricourt », le fait passer la douane et lui
%  fait avancer sa montre de cinquante-cinq minutes. La note du bloc
%  noto-1 le dit deja — « le voyage est decrit tel qu'il etait le long
%  de la frontiere d'avant-guerre ». La colonne allemande garde le nom
%  de la gare tel quel, parce que c'est ainsi qu'elle s'appelait :
%  Avricourt avait ete coupe en deux par la frontiere de 1871, et la
%  moitie allemande s'appelait « Deutsch-Avricourt ». UN NOM DE LIEU
%  QUI EST LUI-MEME UNE FRONTIERE. Les cinquante-cinq minutes sont
%  l'ecart de l'heure de Paris a l'heure d'Europe centrale.
%
%  ET LE VOCABULAIRE DU CHEMIN DE FER EST ALLEMAND, MAIS IL A VOYAGE
%  D'UNE TROISIEME FACON, QUE LA SERIE N'AVAIT PAS ENCORE VUE.
%
%  « Bahnhof » est un compose de 1840 : la cour de la voie. Rien dans
%  la chose ne suggere cette image-la — une gare n'est pas une cour —
%  et pourtant quatre voisins l'ont, chacun avec ses propres racines :
%     hongrois   « pályaudvar »  = voie + cour
%     croate     « kolodvor »    = roue + cour
%     tcheque    « nádraží »     = ce-qui-est-au-bord-de-la-voie
%     slovene    « kolodvor »    idem
%  Aucun n'a pris le mot allemand ; tous ont pris SA FORME. C'est un
%  calque, non un emprunt.
%
%  CELA TRANCHE UNE QUESTION LAISSEE OUVERTE AU TABLEAU 2. La-bas,
%  trois langues sans rapport — l'allemand, l'estonien, le romanche —
%  avaient le meme compose pour le football, pied plus balle, et l'on
%  avait conclu que la chose se decrit d'elle-meme. Ici l'image est
%  ARBITRAIRE : personne ne trouve « cour » tout seul en regardant une
%  gare. Donc, quand quatre langues ont le meme compose arbitraire,
%  elles se sont copiees ; quand elles ont le meme compose evident,
%  elles ne se sont peut-etre pas copiees du tout.
%
%     LA DIFFERENCE ENTRE LA CONVERGENCE ET LE CALQUE SE MESURE A
%     L'EVIDENCE DE L'IMAGE, ET A RIEN D'AUTRE.
%
%  C'est un moyen de verifier les onze tableaux precedents. Le
%  « Löschpapier » du tableau 1 (papier-qui-eteint) est arbitraire, et
%  l'anglais dit « blotting paper », le francais « buvard » : trois
%  images, aucun calque. La « Feuerwehr » du tableau 11 est arbitraire
%  aussi, et l'estonien a fait le meme compose : la, il faudrait
%  chercher.
%
%  ENFIN, TROIS MOTS DE CETTE PAGE SONT FRANCAIS ET N'ONT PAS BOUGE :
%  Waggon, Perron, Billet. Ce sont les
%  trois que le Sprachverein de 1890 a le plus attaques, avec
%  « Fahrkarte », « Bahnsteig » et « Wagen » pour les remplacer. Il a
%  gagne sur les trois EN ALLEMAGNE et perdu sur les trois en Autriche
%  et en Suisse, ou l'on dit encore Perron et Billett. La colonne
%  ecrit Fahrkarte et Bahnsteig ; le partage est encore une fois la
%  matiere de la colonne autrichienne.
% ===================================================================

%%K t12-apar-1 apar
\VUsaut{4.0mm}
\VUtitre{15.0pt}{60}{TAFEL Nr. 12}
\VUsaut{2.4mm}
\VUpk{11.4pt}{40}{Der Bahnhof. --- Die Eisenbahn. --- Die Schiffe.}
\VUsaut{3.4mm}
\textsc{Anmerkung.} --- \textit{Die Fahrpläne sind in jedem Land anders, deshalb
haben wir die Zeiten der}
\VUgras{französischen Tafel} \textit{als Beispiel genommen.}

%%K t12-01-1 p
1. --- Ich habe gerade eine Reise durch Deutschland gemacht. Vor der Abfahrt hatte
ich ein \VUgras{Kursbuch} \textsuperscript{(15)} gekauft, um die Abfahrtszeiten der
Züge zu erfahren. Als ich herausgefunden hatte, dass der günstigste Zug um neun Uhr
abends in Paris abfährt und um ein Uhr dreizehn in Straßburg ankommt, machte ich
mich zur Reise fertig, und am nächsten Tag ließ ich mich zum Bahnhof fahren.

%%K t12-02-1 p
2. --- Als ich in die große Abfahrtshalle kam, hinter einem alten Landpfarrer
\textsuperscript{(2)} mit seiner \VUgras{Reisetasche} \textsuperscript{(3)}, waren
schon ziemlich viele \VUgras{Reisende} \textsuperscript{(1)} da.

%%K t12-02-2 p
Weil ich ein \VUgras{Telegramm} \textsuperscript{(5)} aufgeben wollte, ließ ich mir
von einem \VUgras{Aufsichtsbeamten} \textsuperscript{(6)} das \VUgras{Telegrafenamt}
\textsuperscript{(4)} zeigen.

%%K t12-03-1 p
3. --- Bevor ich in den \VUgras{Wartesaal} \textsuperscript{(7)} ging, holte ich mir
eine \VUgras{Rückfahrkarte} \textsuperscript{(9)}.

%%K t12-04-1 p
4. --- Am \VUgras{Fahrkartenschalter} \textsuperscript{(8)} musste ich nicht lange
warten, weil kaum Leute da waren. Ein \VUgras{Soldat} \textsuperscript{(10)}, der auf
Urlaub fuhr, ließ mich freundlich vor, so dass ich Zeit hatte, dem
\VUgras{Bahnhofsvorsteher} \textsuperscript{(13)} ein paar Fragen zu stellen, der
gerade mit dem \VUgras{Polizeikommissar} \textsuperscript{(12)} und dem
diensthabenden \VUgras{Gendarmen} \textsuperscript{(11)} sprach.

%%K t12-tit-1 sub
\VUpk{10.2pt}{40}{Das Gepäck wird aufgegeben.}
\VUsaut{3.0mm}

%%K t12-05-1 p
5. --- Nachdem ich am \VUgras{Bücherstand} \textsuperscript{(14)} eine Zeitung gekauft
hatte, ließ ich mein \VUgras{Gepäck} \textsuperscript{(17)} wiegen --- ein
\VUgras{Gepäckträger} \textsuperscript{(16)} hatte es eben auf einem
\VUgras{Gepäckkarren} \textsuperscript{(19)} heraufgebracht. Der \VUgras{Beamte}
\textsuperscript{(22)} an der \VUgras{Waage} \textsuperscript{(21)} sagte mir, mein
Koffer wiege fünfunddreißig Kilogramm, das mache ein \VUgras{Übergewicht}
von fünf Kilogramm, für das ich am \VUgras{Gepäckschalter}
\textsuperscript{(23)} einen Zuschlag zahlen musste.

%%K t12-06-1 p
6. --- Weil ich früh dran war, konnte ich dem Treiben der Reisenden im Bahnhof
zusehen. Die einen gaben kleine Taschen an der \VUgras{Gepäckaufbewahrung}
\textsuperscript{(25)} ab oder holten sie dort; andere sahen in die
\VUgras{Fahrpläne} \textsuperscript{(26)}. Ankommende Reisende eilten zum \VUgras{Ausgang}
\textsuperscript{(32)}, die \VUgras{Handkoffer} \textsuperscript{(24)} in der
Hand. Ein paar warteten an der \VUgras{Gepäckausgabe}
\textsuperscript{(27)} und gaben ihren \VUgras{Schein} \textsuperscript{(29)} ab, um
ihre Koffer, \VUgras{Kisten} \textsuperscript{(28)} oder \VUgras{Körbe}
\textsuperscript{(30)} zu bekommen.

%%K t12-07-1 p
7. --- Die \VUgras{Zollbeamten} \textsuperscript{(33)} sahen die Pakete sorgfältig
durch, ob etwas zu verzollen war.

%%K t12-08-1 p
8. --- Manche Reisende gingen in die \VUgras{Bahnhofswirtschaft}
\textsuperscript{(34)}, wo ein \VUgras{Kellner} \textsuperscript{(35)} eifrig
bediente. Sie mussten sich beeilen, weil der \VUgras{Zug} \textsuperscript{(36)}, ein
Schnellzug, nicht lange im Bahnhof halten würde.

%%K t12-09-1 p
9. --- Dann ging ich auf den Abfahrtsbahnsteig hinaus und suchte einen
durchgehenden \VUgras{Wagen} \textsuperscript{(37)}, um unterwegs nicht umsteigen zu
müssen. Ich stieg in einen Durchgangswagen mit einem \VUgras{Abteil}
\textsuperscript{(38)} für Raucher und einem für „allein reisende Damen“.

%%K t12-tit-2 sub
\VUpk{10.2pt}{40}{Abfahrt bei Nacht.}
\VUsaut{3.0mm}

%%K t12-10-1 p
10. --- Nachdem ich das \VUgras{Trittbrett} \textsuperscript{(40)} hinaufgestiegen
war, die \VUgras{Tür} \textsuperscript{(39)} geschlossen, das \VUgras{Rollo}
\textsuperscript{(41)} hochgezogen und meine kleinen Taschen in das
\VUgras{Gepäcknetz} \textsuperscript{(43)} gelegt hatte, setzte ich mich auf die
\VUgras{Bank} \textsuperscript{(42)}. Als ich den \VUgras{Lampenwärter}
\textsuperscript{(44)} die Lampen anzünden sah, fiel mir ein, dass ich eine Nacht im
Wagen zu verbringen hatte, und ich rief den Schaffner, der die \VUgras{Kissen}
\textsuperscript{(45)} vermietet, um ihn um eines zu bitten.

%%K t12-11-1 p
11. --- Der Schaffner kam, um die Fahrkarten zu prüfen. Ich fragte ihn, warum wir
noch nicht abgefahren seien, da es doch schon Zeit war. Er antwortete, der
\VUgras{Zugführer} \textsuperscript{(46)} sei damit beschäftigt, Fahrräder in den
\VUgras{Packwagen} \textsuperscript{(47)} bringen zu lassen. Außerdem seien noch
nicht alle \VUgras{Wärmflaschen} \textsuperscript{(48)} gewechselt.

%%K t12-12-1 p
12. --- Aber wir waren im Begriff abzufahren: der \VUgras{Heizer}
\textsuperscript{(50)}, oben auf seinem \VUgras{Tender} \textsuperscript{(49)}, holte
Kohle, um das Feuer der \VUgras{Lokomotive} \textsuperscript{(54)} zu schüren, und
\VUgras{der Lokführer} \textsuperscript{(51)} wartete nur noch auf das Zeichen des
\VUgras{Vorstehers} \textsuperscript{(52)}, der auf dem \VUgras{Bahnsteig}
\textsuperscript{(53)} stand.

%%K t12-13-1 p
13. --- Der \VUgras{Kessel} \textsuperscript{(56)} der Lokomotive brummte dumpf; der
\VUgras{Schornstein} \textsuperscript{(55)} stieß Ströme von Rauch aus, mit
\VUgras{Dampf} \textsuperscript{(60)} vermischt. Ein schriller \VUgras{Pfiff}
\textsuperscript{(59)} ertönte, und die \VUgras{Kolben} \textsuperscript{(58)} setzten
sich in Bewegung und trieben die \VUgras{Räder} \textsuperscript{(57)} der schweren
Maschine an.

%%K t12-tit-3 sub
\VUsaut{3.0mm}
\VUpk{10.2pt}{40}{Und los!}
\VUsaut{3.0mm}

%%K t12-14-1 p
14. --- Der Zug lief auf den \VUgras{Schienen} \textsuperscript{(64)} und wurde
schneller; die \VUgras{Bahnsteige} \textsuperscript{(65)} hörten auf; wir kamen an den
\VUgras{Aborten} \textsuperscript{(66)} vorbei und auch an einer Reihe von Wagen, die
auf einem anderen \VUgras{Gleis} \textsuperscript{(63)} standen:
\VUgras{Schlafwagen} \textsuperscript{(67)}, ein \VUgras{Speisewagen}
\textsuperscript{(68)}, ein \VUgras{Postwagen} \textsuperscript{(69)}.

%%K t12-noto-1 noto
\VUnotes{143.2mm}{%
\textit{(*) Anmerkung.} --- \textit{Die Reise ist natürlich so beschrieben, wie sie
entlang der Vorkriegsgrenze verlief.}%
}

%%K t12-15-1 p
15. --- Dann zog der \VUgras{Güterbahnhof} \textsuperscript{(71)} an unseren Augen
vorbei. Unter den Hallen luden Beamte die \VUgras{Güter} \textsuperscript{(72)} auf,
die als Frachtgut gingen. \VUgras{Rangierer} \textsuperscript{(76)} neben dem
\VUgras{Wasserturm} \textsuperscript{(73)} schoben \VUgras{Viehwagen}
\textsuperscript{(75)} und \VUgras{Flachwagen} \textsuperscript{(79)} zusammen, um
einen \VUgras{Güterzug} \textsuperscript{(80)} zu bilden.

%%K t12-16-1 p
16. --- Wir fuhren über die letzte \VUgras{Weiche} \textsuperscript{(78)}, neben der
der Wärter stand; wir kamen am \VUgras{Signal} \textsuperscript{(86)} vorbei, und der
Bahnhof verschwand in der Ferne.

%%K t12-17-1 p
17. --- Bald fuhren wir auf die \VUgras{Eisenbrücke} \textsuperscript{(74)} hinaus,
die über den \VUgras{Fluss} \textsuperscript{(81)} führt. Jenseits der Brücke fuhren
wir am Fluss entlang, auf dem kräftige \VUgras{Schleppdampfer}
\textsuperscript{(82)} schwere \VUgras{Kähne} \textsuperscript{(83)} zogen. Wir sahen
auch ein \VUgras{Fahrgastschiff} \textsuperscript{(84)}, das gerade an einer
\VUgras{Landebrücke} \textsuperscript{(85)} anlegte.

%%K t12-18-1 p
18. --- Nachdem wir durch mehrere Bahnhöfe gerast waren, fuhren wir durch ein paar
\VUgras{Tunnel} \textsuperscript{(87)} und ließen unzählige \VUgras{Häuschen}
\textsuperscript{(88)} der \VUgras{Bahnwärter} hinter uns. Vom Schaukeln des Zuges
gewiegt, fiel ich in tiefen Schlaf.

%%K t12-tit-4 sub
\VUsaut{3.0mm}
\VUpk{10.2pt}{40}{Bahnhof Deutsch-Avricourt.}
\VUsaut{3.0mm}

%%K t12-19-1 p
19. --- Ich wurde plötzlich von der Stimme eines Beamten geweckt, der rief:
„Deutsch-Avricourt! \textsuperscript{(*)} Alles aussteigen!“ Es war die
Zollkontrolle mit allen lästigen Förmlichkeiten der Grenze. Ich musste meine Uhr
nach der \VUgras{Bahnhofsuhr} \textsuperscript{(89)} stellen, die
fünfundfünfzig Minuten vorging; kaum wach, konnte ich den \VUgras{großen Zeiger}
\textsuperscript{(90)} kaum vom \VUgras{kleinen} \textsuperscript{(91)}
unterscheiden; das \VUgras{Zifferblatt} \textsuperscript{(92)} selbst war vom
Morgennebel ganz verschwommen.

%%K t12-20-1 p
20. --- Während die Zollbeamten ihre Kontrolle machten, sah ich gähnend die
\VUgras{Plakate} \textsuperscript{(93)} an, die die Bahnhofswände schmücken. Ich
frühstückte, um die Zeit herumzubringen, sah auf der Karte des deutschen
\VUgras{Streckennetzes} \textsuperscript{(94)} nach, wie weit es noch bis Straßburg
war, und stieg wieder in meinen Wagen, ermutigt von der Hoffnung, bald am Ziel zu
sein.
\VUsaut{6.0mm}
\VUfilet{22mm}
